% The M05 pitch code sheet: one page for groups 3 and 4. They type this code
% into the empty cell of pitch-seats.ipynb by hand.
%
% Compile with: xelatex -interaction=nonstopmode pitch-code.tex   (run twice)
%
% The code here is the code that was run: tools/build_m05_pitch_notebook.py
% defines `show`, `check`, `graph` and `graph_tool`, and the lines below were
% executed against them. If one of those names changes, change this sheet too.
\documentclass[a4paper, 11pt]{extarticle}
\usepackage[top=0.6in, bottom=0.5in, left=0.8in, right=0.8in]{geometry}
\usepackage{enumitem}
\usepackage{tcolorbox}
\tcbuselibrary{listings}
\usepackage{fontspec}
\usepackage{xcolor}

\IfFontExistsTF{XCharter-Roman.otf}{
  \setmainfont{XCharter}[
    Extension      = .otf,
    UprightFont    = *-Roman,
    BoldFont       = *-Bold,
    ItalicFont     = *-Italic,
    BoldItalicFont = *-BoldItalic]
}{
  \IfFontExistsTF{Charter}{\setmainfont{Charter}}{
    \IfFontExistsTF{Iowan Old Style}{\setmainfont{Iowan Old Style}}{
      \IfFontExistsTF{Georgia}{\setmainfont{Georgia}}{}}}}
% A code face that tells 0 from O and l from 1, because it is copied by hand.
\IfFontExistsTF{JetBrains Mono}{\setmonofont{JetBrains Mono}[Scale=0.92]}{
  \IfFontExistsTF{DejaVu Sans Mono}{\setmonofont{DejaVu Sans Mono}[Scale=0.9]}{
    \IfFontExistsTF{Menlo}{\setmonofont{Menlo}[Scale=0.92]}{}}}

\setlength{\parindent}{0pt}
\setlength{\parskip}{0.35em}
\setlist{nosep, leftmargin=1.6em, itemsep=0.2em}
\pagestyle{empty}

\newtcblisting{code}{
  listing only, colback=white, colframe=black, boxrule=0.8pt, arc=0pt,
  left=8pt, right=8pt, top=5pt, bottom=5pt,
  listing options={basicstyle=\ttfamily\normalsize, columns=fullflexible, keepspaces=true,
    showstringspaces=false, escapeinside={(*}{*)}}}

\begin{document}

{\LARGE\bfseries Module 5 --- Code for group 3}

\vspace{0.2em}
\rule{\textwidth}{0.8pt}

Type each framed box into the empty cell of your section, by hand, and press Shift+Enter. Typing a line is how you find out what it does. The network is called \texttt{graph}. \texttt{show} and \texttt{check} are already defined.

\medskip
{\Large\bfseries igraph}\par
\vspace{-0.3em}\rule{\textwidth}{0.4pt}

\textbf{Box 1}
\begin{code}
labels = graph.community_leiden(
    objective_function="modularity",
    n_iterations=-1).membership
show(labels)
\end{code}
\begin{itemize}
\item \texttt{community\_leiden} searches for the grouping with the highest modularity. \texttt{.membership} is the list of group numbers, one per node.
\item \texttt{show} draws it and prints the matrix of sizes and edge counts.
\end{itemize}

\textbf{Box 2} --- the same, with Louvain. Run it again with another seed.
\begin{code}
import random
random.seed(1)
labels = graph.community_multilevel().membership
show(labels)
check(labels)
\end{code}
\begin{itemize}
\item Louvain starts from a random order of the nodes, so \texttt{seed} can change the answer.
\item \texttt{check} prints the sizes, the cut values and the modularity $Q$.
\end{itemize}

\medskip

\newpage

{\LARGE\bfseries Module 5 --- Code for group 4}

\vspace{0.2em}
\rule{\textwidth}{0.8pt}

Type each framed box into the empty cell of your section, by hand, and press Shift+Enter. Typing a line is how you find out what it does. The network is called \texttt{graph}. \texttt{show} and \texttt{check} are already defined.

\medskip
{\Large\bfseries graph-tool}\par
\vspace{-0.3em}\rule{\textwidth}{0.4pt}

Run \emph{Step A} first.

\textbf{Box 1}
\begin{code}
def fit(g):
    state = gt.minimize_blockmodel_dl(g)
    return [int(b) for b in state.get_blocks().a]

labels = graph_tool(fit)
show(labels)
\end{code}
\begin{itemize}
\item \texttt{fit} runs inside graph-tool. There, \texttt{g} is the network and \texttt{gt} is graph-tool.
\item \texttt{minimize\_blockmodel\_dl} finds the grouping of the stochastic block model with the shortest description of the network. The number of groups is chosen by the fit.
\item \texttt{state.get\_blocks()} holds the group of each node. \texttt{graph\_tool(fit)} runs \texttt{fit} there and returns its result.
\end{itemize}

\textbf{Box 2} --- change the \texttt{state = \ldots} line of \texttt{fit}, run again, and compare.
\begin{code}
state = gt.minimize_blockmodel_dl(
    g, state_args=dict(deg_corr=True))
state = gt.minimize_blockmodel_dl(
    g, multilevel_mcmc_args=dict(B_min=2, B_max=2))
\end{code}
\begin{itemize}
\item The first lets every node have its own degree. The second forces exactly 2 groups. Use one line at a time.
\end{itemize}

\end{document}
